% Generated draft: requires mathematical and source review.
\documentclass[11pt]{article}
\usepackage[margin=25mm]{geometry}
\usepackage{fontspec}
\setmainfont{Latin Modern Roman}
\usepackage{amsmath,amssymb,mathtools,unicode-math}
\setmathfont{Latin Modern Math}
\usepackage{enumitem,needspace,xurl,hyperref}
\setlength{\parindent}{0pt}
\setlength{\parskip}{5pt}
\setlength{\emergencystretch}{3em}
\newcommand{\N}{\mathbb{N}}
\newcommand{\Z}{\mathbb{Z}}
\newcommand{\Q}{\mathbb{Q}}
\newcommand{\R}{\mathbb{R}}
\newcommand{\C}{\mathbb{C}}
\newcommand{\F}{\mathbb{F}}
\newcommand{\A}{\mathbb{A}}
\newcommand{\PP}{\mathbb{P}}
\newcommand{\E}{\mathbb{E}}
\newcommand{\eps}{\varepsilon}
\newcommand{\dd}{\,\mathrm{d}}
\begin{document}
\section*{UnsolvedMath ID: 30005014}
\textbf{Source problem code:} catalogue entry
\subsection{Definitions and mathematical statement}
The supplied problem asks for the dimension of a planar polynomial spline space in an intermediate degree range. Its notation requires conventions that the supplied excerpt does not fully specify.

Under a conventional real polynomial interpretation, let \(\Delta\) be a finite subdivision of a planar region into polygonal cells, and let \(\Delta_2\) denote its two-dimensional cells. For integers \(r,d\geq 0\), put
\[
\mathcal P_d=\{p\in\R[x,y]:\deg p\leq d\},
\]
where degree means total degree. One possible precise definition is
\[
S_d^r(\Delta)=\left\{(p_\sigma)_{\sigma\in\Delta_2}\in\prod_{\sigma\in\Delta_2}\mathcal P_d:
\left.\partial^\alpha p_\sigma\right|_{\sigma\cap\tau}
=\left.\partial^\alpha p_\tau\right|_{\sigma\cap\tau}
\text{ for all }\sigma,\tau\in\Delta_2\text{ and }|\alpha|\leq r
\right\},
\]
where \(\alpha=(\alpha_1,\alpha_2)\in\Z_{\geq0}^2\), \(|\alpha|=\alpha_1+\alpha_2\), and \(\partial^\alpha=\partial_x^{\alpha_1}\partial_y^{\alpha_2}\). Conditions on empty intersections are vacuous. This convention imposes no vanishing conditions on the exterior boundary.

With these conventions, the requested task is: for every admissible planar subdivision \(\Delta\), every integer \(r\geq1\), and every integer \(d\) satisfying
\[
r+1\leq d\leq3r,
\]
determine \(\dim_{\R}S_d^r(\Delta)\), allowing dependence on the actual geometry of \(\Delta\). For \(r=0\), the displayed degree range is empty.

The definition above is an explicit conditional interpretation, not a verified recovery of the report's conventions. The excerpt does not specify the coefficient field, admissible subdivision class, topology of the underlying region, boundary regularity convention, or additional boundary conditions. In particular, it does not settle whether \(\Delta\) must be a triangulation. These qualifications require resolution before the conditional formulation can be treated as the exact source problem.
\subsection{Short English statement}
Determine the dimension of the space of planar piecewise polynomial splines of degree at most \(d\) and smoothness \(r\), for every integer degree \(r+1\leq d\leq3r\). The source leaves the precise subdivision and spline conventions unspecified; the displayed definition provides a conditional interpretation requiring mathematical review.
\subsection{Sources}
\begin{itemize}
\item Local registry: supplied UnsolvedMath source record, permanent ID \(30005014\), title ``Dimensions of Planar Spline Spaces in Intermediate Degrees,'' source problem code \texttt{catalogue entry}.
\item Supplied report citation: \emph{Mini-Workshop: Interpolation, Approximation, and Algebra (2023)}, Oberwolfach Reports, pp. 383--428, DOI \texttt{10.4171/owr/2022/7}; \url{https://doi.org/10.4171/owr/2022/7}. The record separately gives proposed year 2022; the year discrepancy is unresolved.
\item The original excerpt associates the statement with ``On the dimension of trivariate spline spaces,'' Nelly Villamizar, joint work with Michael DiPasquale. Although that presentation title concerns trivariate splines, the quoted problem explicitly concerns planar \(\Delta\). The record does not identify a separate proposer.
\item Supplied literature reference: Tohaneanu, ``smooth planar r-splines''; \url{https://arxiv.org/abs/math/0608449}.
\item Supplied literature reference: Schenck--Stillman--Yuan, ``new spline bounds''; \url{https://arxiv.org/abs/1909.13399}.
\item Literature assessment dated exactly 2026-08-22: the supplied curation labels the problem partially solved, reports dimension formulas and sharp bounds in several degree ranges, including sufficiently large degree, and says that the exact planar dimension for \(r+1\leq d\leq3r\) is not determined in general. This is supplied dated curation, not an independently verified literature assessment. The original excerpt labels the computation open.
\item The registry labels the statement \texttt{corrected\_verified} and reports comparison with the official Oberwolfach report. That verification claim is supplied provenance and has not been independently established here.
\end{itemize}
\end{document}
